% Clinical Background section. Paste into your own thesis, or \input it.
%
% SCOPE: the clinical account of coronary flow, told as the mechanism by which the topology of the
% coronary tree could be a risk factor -- the question this thesis investigates. How that shapes
% the feature choices (LM/LAD/LCX, curvature over the tortuosity index, bifurcation angle together
% with branch diameters) is argued in thesis/week2/hemodynamics.tex.
%
% Rewritten 2026-09-10: refocused on topology as a risk factor and cut to roughly the length of the
% other clinical background sections. Removed, in git if wanted back: the schultz2023ess paragraph
% (its diameter gradient is largely produced by its fixed boundary conditions, see
% docs_thesis/papers/schultz2023ess.md) and the numeric detail on CT-FFR accuracy and outcomes
% (norgaard2014nxt AUC 0.90 vs 0.81; biswas2026ffrct HR 3.94, 95% CI 2.92-5.31; yang2023target
% 28.3% vs 46.2% catheterization, HR 0.88, 95% CI 0.59-1.30).
% The label sec:geometry sits on the topology subsection because 03_cad.tex already refers to it;
% the section it once pointed at (06_geometry_risk.tex) no longer exists.
%
% CITATION WANTED: the bend-relative-to-width point is the Dean number (Dean 1927); no refs.bib
% entry yet, so it is stated uncited.
% SECOND-HAND: the helical-flow sentence rests on shen2026geometry's Table 7, not the primary
% papers. Cite for framing only.
%
% NOTE for the methods chapter: models/ffr_unet.py and configs/ffr_unet.json do NOT compute FFR.
% "FFR-UNet" is Feature-Fusion-and-Rectification 3D-UNet, a lumen segmentation architecture named
% for the group's downstream application. Never describe it as an FFR estimator.
%
% FIGURE WANTED: flow separating at a bifurcation, with the low-shear region on the
% lateral wall marked, and the same for the inner wall of a bend. This one could be
% ours rather than external -- a schematic drawn in matplotlib, no licence needed.
% Do not add a \ref until the figure exists.

\section{Coronary hemodynamics}
\label{sec:hemodynamics}

Coronary hemodynamics is the study of the flow of blood through the coronary arteries and the
mechanical forces that this flow exerts on the vessel wall.
It matters to this thesis because it is the route by which the shape of the coronary tree could
influence where, and in whom, disease develops.
This section describes the two quantities that dominate its clinical use, wall shear stress and
fractional flow reserve, and then the evidence that the topology of the tree shapes the first of
them.
Topology is used here, as throughout this thesis, in a broad sense: how the tree branches, the
angles and diameters at its bifurcations, and how its vessels bend and twist.

\subsection{Wall shear stress}

Blood flowing along an artery drags on the wall it passes.
That drag is \textbf{wall shear stress (WSS)}.
For steady flow of a fluid of viscosity $\mu$ at flow rate $Q$ through a tube of radius $R$,
%
\begin{equation}
  \tau = \frac{4 \mu Q}{\pi R^{3}},
  \label{eq:wss}
\end{equation}
%
so halving the radius at the same flow raises the drag eightfold \cite{rampidis2022geometry}.
An artery therefore does not feel one uniform drag along its length; it changes wherever the vessel
narrows, widens, bends or divides.

The endothelium senses that drag.
The whole coronary tree is exposed to the same cholesterol and the same blood pressure, yet where
the drag is low, or reverses direction over the cardiac cycle, the endothelial cells turn inflamed
and plaque-prone, and it is there that plaque starts and grows \cite{chatzizisis2007ess}.
High drag does different damage: over a plaque that already narrows the vessel, it is associated
with regression of the fibrous cap that stabilizes it, leaving the plaque more prone to rupture
\cite{bajraktari2021highwss}.

\subsection{Fractional flow reserve}
\label{sec:ffr}

\textbf{Fractional flow reserve (FFR)} measures not how narrow a vessel looks, but whether the
narrowing limits the blood supply to the heart muscle beyond it \cite{knuuti2019esc}.
It is the ratio of the pressure just past the narrowing to the pressure before it, measured with a
wire during catheterization while flow is driven artificially high, because at rest the small
vessels beyond a narrowing widen to compensate and hide its effect \cite{duncker2015regulation}.
\textbf{CT-derived fractional flow reserve (CT-FFR)} computes the same ratio by simulating blood
flow through the lumen reconstructed from a CCTA scan.
It identifies flow-limiting narrowings more accurately than judging how narrow the vessel looks on
the same scan \cite{norgaard2014nxt}, and a low value carries roughly four times the risk of a
cardiac event \cite{biswas2026ffrct}.

\begin{figure}[h]
  \centering
  \includegraphics[width=\textwidth]{content/background/figures/wss_schematic.png}
  \caption{The shape of a vessel decides where its wall feels low shear. (a) In a bend, the
    fast-moving core of the flow is carried to the outer wall, which feels high WSS, while flow
    along the inner wall slows and can recirculate. (b) At a bifurcation, the fastest flow enters
    each branch next to the flow divider, which feels high WSS, while flow along the lateral walls
    separates and recirculates. The orange walls are where Asakura and Karino found almost all
    plaque \cite{asakura1990flow}. Schematic; the velocity profiles are drawn by hand, not
    simulated.}
  \label{fig:wss-sites}
\end{figure}